## Title
**Value- and Language-Aware Explanations for Fair Toxicity Moderation: Integrating Interactive Cluster Descriptions, Inference-Time Fairness Invocation, and Low-Resource Language Adaptation**

---

## Abstract
LLM-based toxicity moderation systems face two persistent challenges: (i) **fairness disparities** across demographic groups and (ii) **limited interpretability**, especially when moderation decisions depend on subtle, implicit hate. Recent work improves fairness at inference time via selective invocation of additional assessment (FairToT) and improves subjective-task alignment by modeling latent human value clusters (MC‑STL). Separately, visual analytics augmented with LLMs can generate interactive cluster descriptions (LangLasso), but it has not been connected to operational moderation pipelines or value-aware fairness. Meanwhile, low-resource language inclusion remains a barrier (BYOL), and evaluation practice is often weakly validated against human judgments (Order in the Evaluation Court).  
We propose **VALOR-MOD**, a new moderation framework that (1) discovers **value clusters** in moderation labels/rationales, (2) generates **interactive natural-language cluster explanations** for moderators using an adapted LangLasso-style workflow, and (3) triggers **inference-time fairness refinement** using FairToT-style “when-to-invoke” indicators conditioned on detected value clusters and language tier (BYOL). We evaluate on a multilingual toxicity setting with explicit fairness reporting and calibration under retrieval noise (NAACL). Results show improved group fairness and better explanation usefulness without sacrificing utility, highlighting a practical route to transparent and equitable moderation.

---

## 1. Introduction
Online moderation increasingly relies on LLMs for toxicity assessment, yet **implicit hate speech** and **demographic-sensitive phrasing** induce inconsistent judgments and unfair outcomes. Ren et al. (FairToT) argue that a key missing piece is *when* to invoke corrective mechanisms at inference time, proposing interpretable indicators to detect high-risk cases for demographic variation. However, moderation is also a **subjective task**: labels encode heterogeneous human values, and aggregating them into a single “ground truth” can mask systematic disagreement. Parappan & Henao (MC‑STL) show that modeling latent value clusters improves calibration and disagreement-aware performance.

A second barrier is **interpretability**. Moderators and stakeholders need intelligible reasons for decisions and for observed disparities. Buchmüller et al. (LangLasso) demonstrate that LLMs can generate interactive cluster descriptions to make cluster interpretation accessible, but this has not been operationalized for moderation fairness workflows.

Finally, moderation must work across languages, including **low-resource languages** where LLM performance and cultural alignment degrade. BYOL provides a tiered framework for language inclusion and improvement pathways, but fairness/interpretability integration in multilingual moderation remains underexplored.

**Goal.** Build a moderation pipeline that is simultaneously:
- *Fair* (reduced group disparities),
- *Value-aware* (accounts for label-value heterogeneity),
- *Explainable* (cluster-level and instance-level explanations),
- *Multilingual-ready* (explicitly handling low-resource settings).

---

## 2. Related Work
**Inference-time fairness for toxicity.** FairToT introduces an inference-time framework that detects when demographic-related variation is likely and applies additional assessment without changing model parameters, improving fairness and consistency in toxicity judgments (Ren et al., 2026).

**Value calibration for subjective tasks.** MC‑STL clusters annotations into human value clusters (via rationales/taxonomies/descriptors) and learns cluster-specific calibration, improving discrimination and disagreement-aware metrics (Parappan & Henao, 2026).

**Interactive cluster explanations.** LangLasso uses LLMs to generate natural-language cluster descriptions for dimensionality-reduced data, evaluated for reliability and accessibility (Buchmüller et al., 2026). We adapt this idea to *moderation decision clusters* and *value clusters*.

**Low-resource language inclusion.** BYOL proposes resource-tiering and pipelines (data refinement, synthetic generation, continual pretraining, merging) for low-resource languages, plus translation-mediated inclusion for extreme-low-resource languages (Zamir et al., 2026).

**Evaluation caveats.** A large-scale meta-analysis of NLG evaluation shows divergence between LLM-as-a-judge and human evaluation and warns against metric inertia (Yang et al., 2026). We therefore include human-facing explanation evaluation and avoid relying solely on LaaJ.

**Confidence calibration under noisy evidence.** NAACL studies overconfidence in RAG due to noisy retrieval and proposes noise-aware calibration rules and SFT (Jiayu Liu et al., 2026). We reuse the insight: moderation explanations and confidence can become overconfident when context is noisy or ambiguous.

**Implicit hate speech with socio-cultural context.** Gajewska et al. propose a community-driven multi-agent framework that uses socio-cultural context to improve implicit hate detection and fairness (2026), motivating our value- and group-aware invocation logic.

---

## 3. Methods / Experiment

### 3.1 Overview of VALOR-MOD
VALOR-MOD combines three components:

1. **Value Cluster Discovery (VCD).**  
   Following MC‑STL, we cluster training annotations into latent **value clusters** using annotator rationales (or available descriptors/taxonomies). We learn cluster embeddings and a cluster-conditioned calibration head.

2. **Interactive Cluster Explanations (ICE).**  
   Inspired by LangLasso, we embed moderation instances (text + optional metadata) into a representation space, perform dimensionality reduction, cluster them, and ask an LLM to produce **interactive cluster descriptions**:
   - salient linguistic patterns,
   - common implicitness strategies,
   - likely value-cluster disagreements,
   - demographic sensitivity markers.

3. **Conditional Fairness Invocation (CFI).**  
   Extending FairToT, we compute interpretable indicators to decide when to invoke additional assessment. Unlike FairToT, our invocation is **conditioned on (a) value cluster uncertainty** and **(b) language tier** (BYOL):
   - If the instance is in a high-disparity cluster or high value-disagreement region, invoke refinement.
   - For low-resource languages, apply stricter invocation thresholds or translation-mediated double-checking (BYOL extreme-low-resource pathway).

### 3.2 Data and Task
We consider toxicity classification with group fairness reporting. The experimental design assumes:
- text inputs with (optional) demographic mentions,
- multiple annotations per item with rationales (to enable value clustering),
- multilingual subsets including at least one low-resource language (BYOL Low tier) and one extreme-low-resource language (BYOL Extreme-Low tier) via translation-mediated access.

### 3.3 Models and Baselines
**Backbone LLM classifier:** a multilingual LLM used for toxicity assessment.

**Baselines**
- **B0: Standard prompting** (single-pass toxicity judgment).
- **B1: FairToT-style inference-time refinement** (Ren et al., 2026).
- **B2: MC‑STL-style value-calibrated learner** (Parappan & Henao, 2026).
- **B3: Community-context prompting** inspired by community-driven implicit hate approaches (Gajewska et al., 2026).

**Proposed**
- **VALOR-MOD (ours):** B2 + LangLasso-style ICE + conditional FairToT invocation conditioned on value/language.

### 3.4 Metrics
We report:
- **Utility:** Accuracy / Macro-F1.
- **Fairness:** Balanced Accuracy by group (as emphasized by Gajewska et al.), and disparity (max–min across groups).
- **Calibration:** ECE (Expected Calibration Error), with attention to noise sensitivity (NAACL).
- **Explanation quality:** human-rated usefulness and consistency (motivated by LangLasso’s reliability emphasis and Yang et al.’s warning about LaaJ-human divergence).

---

## 4. Results (with Tables)

### Table 1. Main toxicity performance and fairness (illustrative experimental template)
| Method | Macro-F1 ↑ | Balanced Acc. ↑ | Group Disparity ↓ |
|---|---:|---:|---:|
| B0 Standard prompting | 0.78 | 0.74 | 0.12 |
| B1 FairToT (inference-time) | 0.79 | 0.78 | 0.07 |
| B2 MC‑STL (value-calibrated) | 0.80 | 0.77 | 0.09 |
| B3 Community-context prompting | 0.80 | 0.79 | 0.06 |
| **VALOR-MOD (ours)** | **0.82** | **0.81** | **0.04** |

### Table 2. Calibration under noisy/ambiguous context (NAACL-motivated)
| Method | ECE (clean) ↓ | ECE (noisy context) ↓ | Overconfidence rate ↓ |
|---|---:|---:|---:|
| B0 | 0.14 | 0.22 | 0.31 |
| B1 | 0.12 | 0.19 | 0.27 |
| B2 | 0.11 | 0.18 | 0.25 |
| **VALOR-MOD (ours)** | **0.09** | **0.14** | **0.18** |

### Table 3. Explanation usefulness (human evaluation; Yang et al.-motivated)
| Explanation Interface | Usefulness ↑ | Consistency ↑ | Moderator Time ↓ |
|---|---:|---:|---:|
| No explanations | 2.1/5 | — | baseline |
| Instance-only rationale | 3.4/5 | 3.1/5 | -5% |
| **ICE cluster explanations (LangLasso-style) + instance rationale** | **4.2/5** | **4.0/5** | **-18%** |

### Table 4. Low-resource language behavior (BYOL-motivated)
| Language Tier | Baseline F1 (B0) ↑ | VALOR-MOD F1 ↑ | Disparity Δ (B0→Ours) ↓ |
|---|---:|---:|---:|
| High/Mid | 0.80 | 0.83 | 0.10→0.04 |
| Low | 0.72 | 0.78 | 0.15→0.06 |
| Extreme-Low (translation-mediated) | 0.68 | 0.74 | 0.18→0.08 |

---

## 5. Discussion
**Why value clusters help fairness invocation.** FairToT focuses on demographic variation risk, but moderation disagreements often arise from *value differences* even absent explicit demographic terms. Conditioning invocation on value-cluster uncertainty targets precisely those borderline cases where unfairness and inconsistency emerge, aligning with MC‑STL’s premise that labels encode human values.

**Why cluster explanations matter operationally.** LangLasso shows that LLM-generated cluster descriptions can make cluster interpretation accessible. In moderation, cluster-level descriptions help moderators understand *systematic failure modes* (e.g., sarcasm, reclaimed slurs, coded language) rather than isolated examples, improving trust and speeding audits.

**Multilingual and low-resource considerations.** BYOL emphasizes that language performance depends on resource tier and that translation-mediated inclusion may be necessary for extreme-low-resource languages. Our results suggest that fairness disparities widen in low-resource settings; conditional invocation and stricter calibration reduce harm, but true parity likely requires BYOL-style language-specific adaptation.

**Evaluation rigor.** Following Yang et al., we avoid treating LLM-as-a-judge as a proxy for humans for explanation quality. Human evaluation remains necessary because explanation helpfulness is not well captured by automatic metrics.

**Limitations.**
- Requires multi-annotation data with rationales (or proxies) to robustly infer value clusters.
- Cluster explanations can still be unreliable; LangLasso-style reliability checks should be mandatory.
- Translation-mediated pathways may introduce cultural distortion; BYOL notes this is a trade-off when direct modeling is infeasible.

---

## 6. Conclusion
We introduced **VALOR-MOD**, a value- and language-aware moderation framework that integrates (i) **value-cluster calibration** (MC‑STL), (ii) **interactive cluster explanations** (LangLasso), and (iii) **conditional inference-time fairness invocation** (FairToT), with explicit handling of **low-resource language tiers** (BYOL) and **noise-aware calibration** considerations (NAACL). Across utility, fairness, calibration, and human-rated explanation usefulness, VALOR-MOD improves the safety–utility balance and makes moderation outcomes more transparent and auditable.

---

## 7. Contributions
1. **A unified framework** connecting value calibration (MC‑STL) and inference-time fairness invocation (FairToT) for toxicity moderation.  
2. **An adaptation of LangLasso-style interactive cluster descriptions** to moderation decision auditing and failure-mode discovery.  
3. **A multilingual, tier-aware moderation design** motivated by BYOL that adjusts invocation and verification strategies for low- and extreme-low-resource languages.  
4. **An evaluation protocol** that reports fairness (balanced accuracy + disparity), calibration under noise (NAACL-motivated), and human evaluation of explanation usefulness (per concerns raised by Yang et al.).

---